\reading{LTR-11}{Contingency Funding Planning}{DEFER}{0}

\begin{lrkeybox}[Learning objectives — official 2026 spine wording]
\begin{enumerate}[label=\textbf{\alph*.}]
\item Discuss the relationship between contingency funding planning and liquidity stress testing.
\item Describe best practices in the design of a sound contingency funding plan.
\item Assess the key components of a contingency funding plan (governance and oversight, scenarios and liquidity gap analysis, contingent actions, monitoring and escalation, and data and reporting).
\end{enumerate}
{\footnotesize\color{FRMGrey}Source: Shyam Venkat and Stephen Baird, \textit{Liquidity Risk Management: A Practitioner's Perspective}, Chapter 7.}
\end{lrkeybox}

\begin{lrtrapbox}[Label collision and stale wording]
The Lecture layer labels this reading \texttt{LO 69.a--c}; the Crash layer uses \texttt{69} for \textbf{LTR-6 Intraday Liquidity Risk Management}. Both layers also render objective \textbf{b} as ``\emph{evaluate the key design considerations}''; the 2026 spine says ``\emph{describe best practices in the design}''. Renumbered and rewarded to the spine.
\end{lrtrapbox}

% =====================================================================
\lo{LTR-11 a}{Discuss the relationship between contingency funding planning and liquidity stress testing.}

\begin{lrdefbox}[The link]
A \term{contingency funding plan (CFP)} serves as a logical connection to its companion, the \textbf{liquidity stress testing framework}, by linking the stress test results and other related information as \textbf{inputs} to the CFP's governance, its menu of contingent liquidity actions, and its decision framework.
\end{lrdefbox}

\begin{center}
\small
\begin{tabularx}{\textwidth}{@{}p{5.3cm} X@{}}
\toprule
\rowcolor{FRMSoft}\textbf{Event profile} & \textbf{Managed by} \\
\midrule
\textbf{Low-impact, high-probability} & Business-as-usual (BAU) funding and liquidity risk management activities. \\
\addlinespace
\textbf{High-impact, low-probability} & The \textbf{contingency funding plan}. \\
\bottomrule
\end{tabularx}
\end{center}
\figcap{Institutions use CFPs to address ``the other end of the spectrum''. The whole relationship sits in this two-row table: the stress test produces the inputs, and the CFP is the instrument for the tail that BAU management does not cover.}

\begin{lrtrapbox}[The impact and probability words are both live]
The CFP handles \textbf{high-impact, low-probability} events; BAU handles \textbf{low-impact, high-probability} ones. Both halves of each pair are corruptible independently, giving four plausible-looking statements of which only one is right. Say both words before locking.
\end{lrtrapbox}

% =====================================================================
\lo{LTR-11 b}{Describe best practices in the design of a sound contingency funding plan.}

\begin{center}
\small
\begin{tabularx}{\textwidth}{@{}p{4.4cm} X@{}}
\toprule
\rowcolor{FRMSoft}\textbf{Design consideration} & \textbf{What it requires} \\
\midrule
\textbf{1. Aligned to business and risk profiles} & The CFP must fit the institution it belongs to. \\
\addlinespace
\textbf{2. Integrated with broader risk management frameworks} & Including \textbf{enterprise risk management (ERM)}, \textbf{capital management}, and \textbf{business continuity and crisis management}. \\
\addlinespace
\textbf{3. Operational and actionable, but a flexible playbook} & Includes a \textbf{menu} of possible contingency actions that management can undertake \emph{in different stress scenarios and at graduated levels of severity}. \\
\addlinespace
\textbf{4. Inclusive of appropriate stakeholder groups} & Various management committees --- ALCO, risk and capital committee, investment committee --- plus business units, finance, treasury, risk, operations and technology. \\
\addlinespace
\textbf{5. Supported by a communication plan} & As in any crisis, the \textbf{coordinated and timely} communication of information to stakeholders is critical. \\
\bottomrule
\end{tabularx}
\end{center}

\begin{lrtrapbox}[``Playbook'' and ``flexible'' are both in the requirement]
Consideration 3 is a two-sided claim: the CFP must be \textbf{operational and actionable} \emph{and} \textbf{flexible}. A distractor keeping only the first half turns the CFP into a fixed script, which is precisely what the reading says it must not be --- the actions are a \emph{menu}, graduated by severity, not a fixed sequence.
\end{lrtrapbox}

% =====================================================================
\lo{LTR-11 c}{Assess the key components of a contingency funding plan (governance and oversight, scenarios and liquidity gap analysis, contingent actions, monitoring and escalation, and data and reporting).}

\subsection{1. Governance and oversight}

\begin{center}
\small
\begin{tabularx}{\textwidth}{@{}p{4.0cm} X@{}}
\toprule
\rowcolor{FRMSoft}\textbf{Stakeholder} & \textbf{Role} \\
\midrule
\textbf{Corporate treasury} & Oversees risk, funding and liquidity in normal circumstances. \textbf{Activates the CFP}, based on economic analysis and stress test results. \\
\addlinespace
\textbf{Liquidity crisis team (LCT)} & The \textbf{central point of contact}; responsible for continuous monitoring of the institution's liquidity profile. Generally \textbf{designs the CFP} and submits it to the senior management group for review and approval. \\
\addlinespace
\textbf{Management committee} & Comprises senior management. \textbf{Manages the LCT during crises}; liaises with the board of directors, analyses liquidity, and \textbf{approves CFP recommendations}. \\
\addlinespace
\textbf{Board of directors} & \textbf{Provides leadership} to the LCT and management committee during crises; offers guidance on CFP implementation. \\
\bottomrule
\end{tabularx}
\end{center}

\begin{lrtrapbox}[Four bodies, four verbs, and they are all swappable]
Treasury \textbf{activates}. The LCT \textbf{designs and monitors}. The management committee \textbf{approves}. The board \textbf{leads and guides}. This is a \S5B.2 role-misassignment field with four candidates for four verbs, which means sixteen possible wrong pairings and only one right set. It is also a confirmed GARP template pair: the practice papers test ``invoke versus coordinate'' on the CFP, each paper testing one verb and offering the other as the trap.
\end{lrtrapbox}

\begin{lrkeybox}[Communication and coordination]
The various business units should be coordinated and interdependent, so as to produce data in a timely manner for decision-making. \textbf{Confidence is critical}, as demonstrated repeatedly in the financial services industry and most recently during the crisis of 2007--08. Communications should be \textbf{centrally managed}, with each channel owned:
\begin{itemize}
\item \textbf{Business units} $\rightarrow$ clients, counterparties
\item \textbf{Corporate treasury} $\rightarrow$ regulators and supervisors, rating agencies, clearing banks
\item \textbf{Investor relations} $\rightarrow$ investors, the analyst community, public media
\item \textbf{Legal and compliance} $\rightarrow$ regulators and supervisors
\end{itemize}
\end{lrkeybox}

\begin{lrtrapbox}[Two functions talk to regulators, and they are not the same two]
\textbf{Corporate treasury} and \textbf{legal and compliance} both face regulators and supervisors. \textbf{Investor relations} faces investors, analysts and the media --- \emph{not} regulators. An option routing media questions through treasury, or regulatory contact through investor relations, is the standard build on this list.
\end{lrtrapbox}

\begin{lrkeybox}[Policies, procedures, and testing]
The CFP's policies and procedures cover: an \textbf{introduction} (overview and purpose, related policies); \textbf{governance} (roles and responsibilities, approval); a \textbf{stress testing and scenarios} overview; \textbf{monitoring and escalation} (liquidity gap analysis, contingent action); and \textbf{reporting} (including frequency).

On a periodic basis, institutions should evaluate their CFP's \textbf{operational readiness and effectiveness}. Activities that create liquidity --- issuing debt, selling securities from the investment portfolio --- could be tested.
\end{lrkeybox}

\subsection{2. Scenarios and liquidity gap analysis}

CFP scenarios should \textbf{match liquidity stress tests} and connect to \textbf{recovery plans}. Liquidity stress tests cover systemic, idiosyncratic, market and funding risks, across short- and long-term stress periods. CFPs inform management about crises and solutions, drawing from stress test results; they may also explore \textbf{additional scenarios} to enhance contingency plans for less likely liquidity challenges.

\begin{lrtrapbox}[The CFP's scenario set is a superset, not a copy]
CFP scenarios \emph{match} the stress tests \textbf{and may go beyond them}. A statement that the CFP uses exactly the stress test scenarios has dropped the extension; a statement that it uses a wholly separate set has broken the link.
\end{lrtrapbox}

\subsection{3. Contingent actions}

\begin{itemize}
\item Contingent actions in response to a liquidity gap analysis should align with the \textbf{shortfall, its timing, and the associated inflow}. Examples: maintaining credit lines, adjusting lending practices, attracting deposits with higher rates, disposing of assets or business units, reducing lending, issuing debt, drawing on credit lines.
\item The choice depends on the \textbf{nature and magnitude} of the stressed event, with market closures, reduced funding access and deposit withdrawals affecting feasibility.
\item Firms must consider the \textbf{message their actions send to stakeholders}. Some actions, while providing short-term relief, \textbf{might convey financial weakness} and hinder the effectiveness of future contingent actions. Firms should balance immediate liquidity needs against long-term considerations and, where possible, \textbf{defer some actions until later stages of the crisis} to avoid exacerbating liquidity issues.
\end{itemize}

\begin{lrtrapbox}[The signalling point is the examinable one]
The non-obvious content here is that an available action may be deliberately \textbf{held back}, because using it early signals weakness and burns its future effectiveness. An option asserting that a firm should deploy all available contingent actions as early as possible is the intuitive-but-wrong answer, and it is the shape \S10.7 warns about: intuitive, and unsupported by the reading.
\end{lrtrapbox}

\subsection{4. Monitoring and escalation}

Monitor \textbf{both} the current liquidity profile \textbf{and} the anticipated effects of potential liquidity events. Monitoring and escalation involve tracking and measuring liquidity risk through \term{early warning indicators (EWIs)}. Monitoring EWIs helps firms anticipate liquidity problems and act in time; a balanced approach is needed to avoid sending negative signals to stakeholders.

\begin{center}
\small
\begin{tabularx}{\textwidth}{@{}p{4.6cm} X@{}}
\toprule
\rowcolor{FRMSoft}\textbf{Measure family} & \textbf{Content} \\
\midrule
\textbf{Market and business measures} \newline \emph{(external and internal)} &
Macro- and microeconomic indicators, covering a wide range of factors from stock market declines to credit downgrades and deposit run-offs. Industry and competitor analysis can provide additional insight. \\
\addlinespace
\textbf{Liquidity health measures} \newline \emph{(internal)} &
Focus directly on the firm's current and future liquidity strength: projected funding requirements, non-core funding, funding source concentration, short-term liabilities to total assets, overnight borrowings to total assets. \\
\bottomrule
\end{tabularx}
\end{center}

\begin{lrkeybox}[The three escalation levels]
\begin{enumerate}
\item \textbf{Level 1.} Begins with \textbf{increased management oversight}, prompted by issues such as liquidity reductions. Focus is on \emph{future} liquidity measures and \emph{market perception}. Internal and external communication is vital.
\item \textbf{Level 2.} Crises have a \textbf{noticeable negative impact}. Detailed analysis of current liquidity and of competitor responses is conducted. Emphasis shifts to \textbf{immediate actions} for recovery and business maintenance, including liquidity-enhancing measures.
\item \textbf{Level 3.} The firm is focused on \textbf{survival as a going concern}. Significant liquidity-boosting actions have \emph{already} been taken. Reporting requirements become more stringent as the crisis escalates.
\end{enumerate}
\textbf{Movement between levels requires LCT approval.}
\end{lrkeybox}

\begin{lrtrapbox}[Three levels, three verbs, and one gatekeeper]
Level 1 \textbf{watches} (oversight, perception, future measures). Level 2 \textbf{acts} (immediate recovery measures). Level 3 \textbf{survives} (the big actions are already behind it). The tell for Level 3 is the past tense --- ``have already been taken''. And the gatekeeper for every transition is the \textbf{LCT}, not the board and not the management committee.
\end{lrtrapbox}

\subsection{5. Data and reporting}

\begin{itemize}
\item Timely and accurate data on the firm's liquidity position is essential. \textbf{Regular liquidity reporting, possibly daily}, is desirable in some firms, encompassing both quantitative data (liquidity ratios) and qualitative information.
\item Reports should outline the methods for determining \textbf{liquidity coverage for future obligations and cash flows}, considering \textbf{intraday positions}, \textbf{contingent liabilities}, and \textbf{funding source usage}.
\item The \textbf{frequency of reporting can vary depending on the severity of the crisis}.
\end{itemize}

\begin{lrtrapboxnb}[Consolidated trap summary — LTR-11]
\begin{itemize}
\item \textbf{Polarity.} CFP $=$ \textbf{high-impact, low-probability}. BAU $=$ low-impact, high-probability. Both words in each pair are corruptible.
\item \textbf{Role.} Treasury \textbf{activates} the CFP; the LCT \textbf{designs and monitors}, and \textbf{approves movement between escalation levels}; the management committee \textbf{approves} CFP recommendations; the board \textbf{leads and guides}. Confirmed GARP template territory (``invoke versus coordinate'').
\item \textbf{Role.} Investor relations faces investors, analysts and media --- \emph{not} regulators. Treasury and legal/compliance face regulators.
\item \textbf{Scope.} CFP scenarios match the stress tests \textbf{and may extend beyond them}.
\item \textbf{Scope.} Some contingent actions are \textbf{deliberately deferred}, because early use signals weakness. ``Deploy everything immediately'' is the intuitive wrong answer.
\item \textbf{Sibling.} Market and business measures are \emph{external and internal} macro/micro indicators; liquidity health measures are \emph{internal} balance-sheet ratios.
\item \textbf{Definition.} Level 3 is identified by the \emph{past tense} --- the significant liquidity-boosting actions have already been taken.
\item \textbf{Scope.} CFP integration reaches ERM, capital management, \textbf{and business continuity and crisis management}.
\end{itemize}
\end{lrtrapboxnb}
